\documentclass[10pt,a4paper]{amsart}
\usepackage{amsmath,amssymb,amsthm,graphicx}
\usepackage{verbatim}
\newenvironment{reference}{\comment}{\endcomment}
%\newenvironment{reference}{}{}
\newenvironment{slogan}{\comment}{\endcomment}
\newenvironment{history}{\comment}{\endcomment}

% For commutative diagrams we use Xy-pic
\usepackage[all]{xy}

% We use 2cell for 2-commutative diagrams.
\xyoption{2cell}
% Theorem environments.
%
\theoremstyle{plain}
\newtheorem{theorem}[subsection]{Theorem}
\newtheorem{proposition}[subsection]{Proposition}
\newtheorem{lemma}[subsection]{Lemma}

\theoremstyle{definition}
\newtheorem{definition}[subsection]{Definition}
\newtheorem{example}[subsection]{Example}
\newtheorem{exercise}[subsection]{Exercise}
\newtheorem{situation}[subsection]{Situation}

\theoremstyle{remark}
\newtheorem{remark}[subsection]{Remark}
\newtheorem{remarks}[subsection]{Remarks}

\numberwithin{equation}{subsection}

% Macros
%
\def\lim{\mathop{\mathrm{lim}}\nolimits}
\def\colim{\mathop{\mathrm{colim}}\nolimits}
\def\Spec{\mathop{\mathrm{Spec}}}
\def\Hom{\mathop{\mathrm{Hom}}\nolimits}
\def\Ext{\mathop{\mathrm{Ext}}\nolimits}
\def\SheafHom{\mathop{\mathcal{H}\!\mathit{om}}\nolimits}
\def\SheafExt{\mathop{\mathcal{E}\!\mathit{xt}}\nolimits}
\def\Sch{\mathit{Sch}}
\def\Mor{\mathop{\mathrm{Mor}}\nolimits}
\def\Ob{\mathop{\mathrm{Ob}}\nolimits}
\def\Sh{\mathop{\mathit{Sh}}\nolimits}
\def\NL{\mathop{N\!L}\nolimits}
\def\CH{\mathop{\mathrm{CH}}\nolimits}
\def\proetale{{pro\text{-}\acute{e}tale}}
\def\etale{{\acute{e}tale}}
\def\QCoh{\mathit{QCoh}}
\def\Ker{\mathop{\mathrm{Ker}}}
\def\Im{\mathop{\mathrm{Im}}}
\def\Coker{\mathop{\mathrm{Coker}}}
\def\Coim{\mathop{\mathrm{Coim}}}

% Boxtimes
%
\DeclareMathSymbol{\boxtimes}{\mathbin}{AMSa}{"02}

%
% Macros for moduli stacks/spaces
%
\def\QCohstack{\mathcal{QC}\!\mathit{oh}}
\def\Cohstack{\mathcal{C}\!\mathit{oh}}
\def\Spacesstack{\mathcal{S}\!\mathit{paces}}
\def\Quotfunctor{\mathrm{Quot}}
\def\Hilbfunctor{\mathrm{Hilb}}
\def\Curvesstack{\mathcal{C}\!\mathit{urves}}
\def\Polarizedstack{\mathcal{P}\!\mathit{olarized}}
\def\Complexesstack{\mathcal{C}\!\mathit{omplexes}}
% \Pic is the operator that assigns to X its picard group, usage \Pic(X)
% \Picardstack_{X/B} denotes the Picard stack of X over B
% \Picardfunctor_{X/B} denotes the Picard functor of X over B
\def\Pic{\mathop{\mathrm{Pic}}\nolimits}
\def\Picardstack{\mathcal{P}\!\mathit{ic}}
\def\Picardfunctor{\mathrm{Pic}}
\def\Deformationcategory{\mathcal{D}\!\mathit{ef}}
\usepackage{cleveref}
\crefname{originallemma}{original lemma}{original lemmas}
\Crefname{originallemma}{Original lemma}{Original lemmas}
\makeatletter
\newlabel{lemma-quasi-compact-permanence}{{03KS}{source}{}{sourcepointer-1}{}}
\newlabel{lemma-quasi-compact-permanence@cref}{{[originallemma][0][]03KS}{[1][0][]source}{}{}{}}
\makeatother
\crefname{originallemma}{original lemma}{original lemmas}
\Crefname{originallemma}{Original lemma}{Original lemmas}
\makeatletter
\newlabel{lemma-quasi-compact-scheme-theoretic-image}{{082Z}{source}{}{sourcepointer-2}{}}
\newlabel{lemma-quasi-compact-scheme-theoretic-image@cref}{{[originallemma][0][]082Z}{[1][0][]source}{}{}{}}
\makeatother
\crefname{originallemma}{original lemma}{original lemmas}
\Crefname{originallemma}{Original lemma}{Original lemmas}
\makeatletter
\newlabel{lemma-section-immersion}{{03KP}{source}{}{sourcepointer-3}{}}
\newlabel{lemma-section-immersion@cref}{{[originallemma][0][]03KP}{[1][0][]source}{}{}{}}
\makeatother
\title{The scheme theoretic image of a partial section is an isomorphism over the domain of the section}
\author{The Stacks Project Authors}
\date{stacks-project:spaces-morphisms.tex:local-snapshot-20261012; GNU FDL 1.2 or later}
\begin{document}
\maketitle
\setcounter{subsection}{1}
\noindent\textit{Editorial note.} All delivered mathematical text is the author's own, copied from the pinned Stacks Project chapter file spaces-morphisms.tex; no statement or proof step was rewritten, reordered or removed.\par\medskip
\noindent\textit{Editorial note.} The proof uses the interaction of two same-chapter statements (commutation with opens and closedness of sections) to conclude an isomorphism; every step is written out.\par\medskip
\noindent\textit{Editorial note.} The three references of the counted spans are recorded as non-load-bearing navigation pointers; the argument that consumes them stays inside the counted proof span.\par\medskip
\noindent\textit{Editorial note.} Preamble: the three native excerpt spans of sources/stacks-spaces-morphisms/preamble.tex (lines 8-13, 14-18 and 154-217). The chapter file itself carries no preamble text.\par\medskip
\noindent\textit{Allowed standard prerequisites (editorial).}\begin{itemize}
\item Permanence of quasi-compactness for the section s. Lemma lemma-quasi-compact-permanence, used at proof lines 3393-3394. Makes s quasi-compact so that the scheme theoretic image behaves well on opens.
\item The scheme theoretic image commutes with restriction to opens for quasi-compact morphisms. Lemma lemma-quasi-compact-scheme-theoretic-image, used at proof lines 3395-3397. Identifies Y' intersect f\textasciicircum{}\{-1\}(V) with the image of the restricted section.
\item A section of a separated morphism is a closed immersion. Lemma lemma-section-immersion, used at proof lines 3398-3402. Turns the section of f\textasciicircum{}\{-1\}(V) -> V into a closed immersion.
\item The scheme theoretic image of a closed immersion is the closed subspace itself. used at proof lines 3402-3404. Concludes that the image map over V is an isomorphism.
\end{itemize}\medskip
\ifdefined\hypertarget\hypertarget{sourcepointer-1}{}\fi
\noindent\textit{Source reference.} Original lemma 03KS: focused-seeds/AGP/sources/stacks-snapshot-20261011/tags.txt, source lines 4688--4688. Navigation pointer into the Stacks Project; the standard fact it supplies is applied inside the counted proof span and the argument resting on it stays there.\par\medskip
\ifdefined\hypertarget\hypertarget{sourcepointer-2}{}\fi
\noindent\textit{Source reference.} Original lemma 082Z: focused-seeds/AGP/sources/stacks-snapshot-20261011/tags.txt, source lines 10269--10269. Navigation pointer into the Stacks Project; the standard fact it supplies is applied inside the counted proof span and the argument resting on it stays there.\par\medskip
\ifdefined\hypertarget\hypertarget{sourcepointer-3}{}\fi
\noindent\textit{Source reference.} Original lemma 03KP: focused-seeds/AGP/sources/stacks-snapshot-20261011/tags.txt, source lines 4685--4685. Navigation pointer into the Stacks Project; the standard fact it supplies is applied inside the counted proof span and the argument resting on it stays there.\par\medskip
\begin{lemma}
\label{lemma-scheme-theoretic-image-of-partial-section}
Let $S$ be a scheme.
Let $f : X \to Y$ be a separated morphism of algebraic spaces over $S$.
Let $V \subset Y$ be an open subspace such that $V \to Y$ is quasi-compact.
Let $s : V \to X$ be a morphism such that $f \circ s = \text{id}_V$.
Let $Y'$ be the scheme theoretic image of $s$.
Then $Y' \to Y$ is an isomorphism over $V$.
\end{lemma}

\begin{proof}
By Lemma \ref{lemma-quasi-compact-permanence}
the morphism $s : V \to X$ is quasi-compact.
Hence the construction of the scheme theoretic image $Y'$
of $s$ commutes with restriction to opens by
Lemma \ref{lemma-quasi-compact-scheme-theoretic-image}.
In particular, we see that $Y' \cap f^{-1}(V)$ is the
scheme theoretic image of a section of the separated
morphism $f^{-1}(V) \to V$. Since a section of a separated
morphism is a closed immersion
(Lemma \ref{lemma-section-immersion}),
we conclude that
$Y' \cap f^{-1}(V) \to V$ is an isomorphism as desired.
\end{proof}

\end{document}
